\chapter{The basics of Banach manifolds}\label{appendix_Banach_manifolds}

The aim of this appendix is to introduce the basics of infinite-dimensional manifolds. This will be used in chapter \ref{Chapter_inf_dim_sympl_red} to give an infinite-dimensional version of the Marsden-Weinstein theorem. This is in turn useful because it provides a rigorous framework to deal with the moment map picture of the Hitchin-Kobayashi correspondence, which is introduced in chapter \ref{chapter_HYM_equation}. Some results in this appendix have proofs that are analogous to those for finite-dimensional manifolds, and we shall omit them. We refer the reader to \parencite{SergeLang}, which is a book that develops the theory of smooth manifolds in a very general setting, giving most proofs and definitions for Banach manifolds, and particularizing results to finite-dimensional ones when there are substantial differences.

\section{Basic definitions}
 Banach manifolds are manifolds that are locally modelled on Banach spaces. To make sense of smoothness, we first need derivatives on Banach spaces. These are Fréchet derivatives. Let \(\mathbf{E},\mathbf{F}\) be two Banach spaces, we write \(\mathcal{L}(\mathbf{E},\mathbf{F})\) for the Banach space of linear maps defined from \(\mathbf{E}\) to \(\mathbf{F}\), which is a Banach space when endowed with the operator norm. We define \(\mathrm{Aut}(\mathbf{E})\) the space of \textit{toplinear isomorphisms}\index{Toplinear isomorphisms}, that is, vector space isomorphisms of \(\mathbf{E}\) that are continuous and have continuous inverses.

\Definition{\textbf{Fréchet derivative}}{Let \((\mathbf{E}, ||\cdot ||_{\mathbf{E}}), (\mathbf{F}, ||\cdot||_{\mathbf{F}})\) be two Banach spaces, \(U \subset \mathbf{E}\) an open set, and \(f\colon U \to \mathbf{F}\) a continuous map. We say that \(f\) is \textit{Fréchet differentiable}\index{Fréchet differentiable map} at \(x \in U\) if there is a linear map \(Df(x)\in \mathcal{L}(\mathbf{E},\mathbf{F})\) such that
\[
\lim_{h\to 0} \frac{||f(x+h)-f(x)-Df(x)(h)||_{\mathbf{F}}}{||h||_{\mathbf{E}}}=0.
\]
}
Observe that this definition is just the direct generalization to a Banach space of the usual differential of a map; that is, a continuous map is Fréchet differentiable if it has a ``best linear approximation''. Observe that if \(f\) is Fréchet differentiable at every point \(x \in U\), then we may define a map \(Df\colon U \to \mathcal{L}(\mathbf{E},\mathbf{F})\). If this map is continuous, we say that \(f\) is a Fréchet-\(\mathscr{C}^1\) map. We now have a continuous map \(Df\colon U \to F'\), where \(\mathbf{F}'=\mathcal{L}(\mathbf{E},\mathbf{F})\), so we may ask whether this new map is Fréchet differentiable. If it is, we define \(D^2f:=D(Df)\), and if this new map is continuous, we say that \(f\) is a Fréchet-\(\mathscr{C}^2\) map. Inductively, we define \(D^kf:=D^{k-1}(Df), \ k \in \N\), and we say that \(f\) is Fréchet-smooth or a Fréchet-\(\mathscr{C}^\infty\) map if it is a Fréchet-\(\mathscr{C}^k\) map for every \(k \in \N\). Whenever it is clear from the context that we refer to Fréchet regularity, we will omit the word Fréchet. Observe that this reduces to the usual definition of differential if the Banach spaces are eucliean space. If a Fréchet derivative exists, then it is unique.

\Definition{\textbf{Banach manifold}}{Let \(\mathcal{M}\) be a Hausdorff topological space. A \(\mathscr{C}^k\) \textit{atlas}\index{Atlas!infinite dimensional} on \(\mathcal{M}\) is a collection of triplets \(\corch{(U_\alpha, \varphi_\alpha, \mathbf{E}_\alpha)}_{\alpha \in A}\), where each \(U_\alpha\) is an open set, \((\mathbf{E}_\alpha, ||\cdot ||_{\alpha})\) is a Banach space, and \(\varphi_\alpha\colon U_\alpha \to \mathbf{E}_\alpha\) is a continuous map satisfying:
\begin{enumerate}
    \item \(\corch{U_\alpha}_{\alpha \in A}\) is an open covering of \(\mathcal{M}\)
    \item Each \(\varphi_\alpha\colon U_\alpha \to \varphi_\alpha(U_\alpha) \subset  \mathbf{E}_\alpha\) is a homeomorphism onto its image, which is open
    \item If \(U_\alpha \cap U_\beta \neq \emptyset\), then
    \[
    \varphi_\alpha \circ \varphi_\beta^{-1}\colon \varphi_\beta(U_\alpha \cap U_\beta)\to \varphi_\alpha(U_\alpha \cap U_\beta)
    \]
    is a Fréchet \(\mathscr{C}^k\) map.
\end{enumerate}
}
Observe that we do not require the Banach spaces that model the manifold locally to be the same. However, whenever \(U_\alpha \cap U_\beta \neq \emptyset\), it follows from the definition that \(\mathbf{E}_\alpha\) and \(\mathbf{E}_\beta\) are isomorphic (indeed, \(D(\varphi_\alpha \circ \varphi_\beta^{-1})\) is a linear isomorphism between them). When all the spaces are the same (call it \(\mathbf{E}\)), we say that \(\mathcal{M}\) is an \textit{\(\mathbf{E}\)-manifold}. Each connected component of a Banach manifold is an \(\mathbf{E}\)-manifold, and usual finite-dimensional smooth manifolds are \(\R^n\)-manifolds. As in the finite-dimensional case, every atlas is contained within a maximal one. Obvious modifications to the previous definition allow us to define \textit{Hilbert manifolds} and \textit{locally convex manifolds} for manifolds that are locally modelled on those spaces. Observe that Hilbert manifolds are just a special case of Banach manifolds, but this is not the case for locally convex manifolds because we do not have norms to define Fréchet derivatives there. Smooth maps are defined analogously: if \(M\) and \(N\) are two Banach manifolds, and \(f \colon \mathcal{M} \to \mathcal{N}\) is a map, we say that it is smooth if whenever \((U, \varphi)\) and \((V, \psi)\) are smooth charts satisfying \(f(U) \subset V\), then the \textit{localization of \(f\)}\index{Localization of a map} defined by \(\psi\circ f \circ \varphi^{-1}\colon \varphi(U) \to \psi(V)\) is a \(\mathscr{C}^\infty\) map (in the Fréchet sense).


\subsection{The tangent space to a Banach manifold} \label{section_tangent_space_Banach}
In the category of finite-dimensional smooth manifolds, the tangent space to a manifold at a point can be defined as the space of derivations acting on germs of functions at that point, and as the set of equivalence classes of curves passing through that point. In the infinite-dimensional setting, these two definitions are no longer equivalent, and the definition chosen is the second one. With this idea in mind, let \(I \subset \R\) be an interval, and \(\mathcal{M}\) be a Banach manifold, a smooth curve is a smooth map \(\alpha \colon I \to \mathcal{M}\). Let \(\alpha\colon I \to \mathcal{M}\) and \(\beta \colon I' \to \mathcal{M}\) be two smooth curves satisfying \(\alpha(0)=\beta(0)=m \in \mathcal{M}\). We say that they are \textit{equivalent}\index{Curves!equivalent} if there is a chart \((U, \psi, \mathbf{E})\) such that \((\psi \circ \alpha)'(0)=(\psi \circ \beta)'(0)\), where we defined
\begin{equation} \label{def_rel_equiv}
(\psi \circ \alpha)'(t)=[D(\psi \circ \alpha)(t)](1).
\end{equation}
In that case, \eqref{def_rel_equiv} holds for every chart. Observe that the right-hand-side makes sense; \(D(\psi \circ \alpha)(t) \in \mathcal{L}(\R, \mathbf{E})\), so when evaluating at \(1\), the result is a vector in \(\mathbf{E}\). This defines an equivalence relation on the set of all smooth curves passing through \(m\) at $t=0$. A \textit{tangent vector}\index{Tangent!vector} is thus an equivalence relation of curves. We denote it by \([\alpha, m]\), or simply \([\alpha]\) if the point is clear from the context. The set of all tangent vectors is called the \textit{tangent space}\index{Tangent!space} of \(\mathcal{M}\) at \(m\), and denoted \(T_m\mathcal{M}\) as usual. This is indeed a vector space. To endow it with a linear structure, we can use any chart \((U, \phi, \mathbf{E})\), which defines a bijection \(\bar \phi\colon  T_m\mathcal{M} \to \mathbf{E}, [\alpha] \mapsto (\phi \circ \alpha)'(0)\). That is the content of the following lemma.

\Lemma{}{Let \(\mathcal{M}\) be a Banach manifold. \label{Lemma_tang_space_to_Banach_manifold}
\begin{enumerate}
    \item Let \(\phi\colon \mathcal{M} \to (\mathbf{E}, ||\cdot||_{\mathbf{E}})\) be a chart around \(m \in \mathcal{M}\), and set \(p_m:=\phi(m) \in \mathbf{E}\). Then, the following maps are inverses
    \[
    \begin{array}{cccc}
        \bar{\phi}\colon & T_m\mathcal{M} & \longrightarrow & \mathbf{E} \\
        & [\alpha] & \mapsto & (\phi \circ \alpha)'(0)
    \end{array}
    \hspace{2 cm}
    \begin{array}{cccc}
        \bar{\phi}^{-1} \colon & \mathbf{E} & \longrightarrow & T_m\mathcal{M} \\
        & v & \mapsto & [t \mapsto \phi^{-1}(p_m+t\cdot v)].
    \end{array}
    \]
    \item For all charts \(\phi, \psi\) around \(m\), we have \(\bar{\psi} \circ \bar{\phi}^{-1}=D(\psi \circ \phi^{-1})(p_m)\).
    \item \(T_m\mathcal{M}\) admits a unique vector space structure and topology making each \(\bar\phi\) into a linear isomorphism.
\end{enumerate}
}{\parskip=0pt
\begin{proof}
    The proof on \parencite[Lemma 1.41]{Schmeding_2022} can be adapted to this situation.
\end{proof}

\begin{obs}
    As happens in the finite-dimensional case, if \(\mathbf{E}\) is a Banach space, then \(\mathrm{id}_{\mathbf{E}}\colon \mathbf{E} \to \mathbf{E}\) defines a chart, and \(\overline{\mathrm{id}_\mathbf{E}}\) allows us to identify \(\mathbf{E}\) and \(T_e\mathbf{E}\) for every \(e \in \mathbf{E}\). In particular, if \(\mathbf{E}=\R\) we define \(\frac{\diff}{\diff t}|_{t}:=(\overline{\mathrm{id}_\R})^{-1}(1)\in T_t\R\).
\end{obs}}
We can then define the differential of a smooth map; if \(f \colon \mathcal{M} \to \mathcal{N}\) is a smooth map between smooth Banach manifolds, then we define \(\diff_mf\colon T_m\mathcal{M}\to T_{f(m)}\mathcal{N}, [\alpha, m] \mapsto [f \circ \alpha, f(m)]\). This map is linear. Using the identification above, the differential of a map whose codomain is a Banach space is a linear map of the tangent space taking values in the same Banach space. 

Many of the usual constructions of finite-dimensional manifold theory carry to the Banach manifold setting. Firstly, the tangent bundle associated to a Banach manifold is defined as usual:
\[
T\mathcal{M}:=\bigsqcup_{m\in \mathcal{M}} T_m\mathcal{M},
\]
and it is a Banach manifold, endowed with a smooth projection \(\pi([\alpha, m])=m\), whose Banach manifold structure is obtained by requiring that for every chart \((U, \phi, \mathbf{E})\), the map \((\pi^{-1}(U), \Phi)\) is a diffeomorphism, where \(\Phi\colon \pi^{-1}(U)\to \phi(U) \times \mathbf{E}, u \mapsto(\pi(u), \bar\phi(u))\). These maps define an atlas (see \parencite[Sec. 1.1.5]{Schmeding_2022}). Now, vector fields are defined as smooth sections of \(T\mathcal{M}\), and the vector space of all vector fields in \(\mathcal{M}\) is denoted \(\mathfrak{X}(\mathcal{M})\). {\parskip=0pt
\begin{obs} \label{obs_vector_fields_and_differentials_inf_dim}
    We still want to have vector fields acting on smooth real-valued functions to provide a Lie algebra structure on \(\mathfrak{X}(\mathcal{M})\). This requires a few comments.
    \begin{enumerate}
        \item Let \(f \colon \mathcal{M} \to \R\) be a smooth map, then, \(\diff_mf([\alpha, m]) = [f \circ \alpha, f(m)] \in T_{f(m)}\R \cong \R\) where the last isomorphism is given by \(\overline{\mathrm{id}_\R}\). Thus, to see it as a real number, we understand \(\diff_mf([\alpha, m]):=\overline{\mathrm{id}_\R}([f\circ \alpha, f(m)])=(f \circ \alpha)'(0)\), seeing the latter as the derivative at \(0\) of a smooth real function. That is, we essentially recover the standard finite-dimensional result: the differential of a smooth map acting on a vector is the derivative at \(0\) of \(f \circ \alpha\) where \(\alpha\) is a smooth curve satisfying \(\alpha(0)=0\) and \(\alpha'(0)\) is the vector.
        \item A vector field does induce a derivation on the algebra of smooth real-valued functions, which is given by \(X(f)(m):=\diff_mf(X_m)\). 
    \end{enumerate}
\end{obs}
As in the finite-dimensional setting, given two vector fields \(X,Y \in \mathfrak{X}(\mathcal{M})\), there is a unique vector \([X,Y]\) whose action on a smooth function is given by \([X,Y](f):=X(Y(f))-Y(X(f))\) \parencite[Prop. V-1.1 and Prop V-1.3]{SergeLang}. The vector field \([X,Y]\) is called the \textit{Lie bracket} of \(X\) and \(Y\). It is bilinear, antisymmetric, and satisfies the Jacobi identity, and hence defines a Lie algebra structure on \(\mathfrak{X}(\mathcal{M})\) (see \parencite[Cor. 1.4]{SergeLang}). }

\subsection{Immersions, submersions and embeddings}
The notions of immersion, submersion and embedding have generalizations to Banach manifolds. First, let \(\mathcal{M}\) be a Banach manifold, and \(\mathcal{N} \subset \mathcal{M}\). We say that \(\mathcal{N}\) is a \textit{Banach-submanifold}\index{Banach!submanifold} if for every \(m \in \mathcal{N}\), there is a chart \((U, \varphi, \mathbf{E})\) of \(\mathcal{M}\), such that \(m \in \mathcal{M}\), and a closed subspace \(\mathbf{F} \subset \mathbf{E}\) such that 
\[
\varphi(U \cap \mathcal{N})=\varphi(U) \cap \mathbf{F}.
\]
Moreover, given a Banach-space \((\mathbf{E}, ||\cdot ||_{\mathbf{E}})\), and a closed vector subspace \(\mathbf{F} \subset \mathbf{E}\), we say that \(\mathbf{F}\) \textit{splits} in \(\mathbf{E}\), if there is another closed vector subspace \(\mathbf{G}\subset \mathbf{E}\) such that \(\mathbf{E}=\mathbf{F}\oplus \mathbf{G}\). We say that \(\mathbf{G}\) is a \textit{topological complement}\index{Topological complement} of \(\mathbf{F}\) in \(\mathbf{E}\). Observe that every closed subspace of a Hilbert space splits (just take its orthogonal complement).

\Definition{}{Let \(\mathcal{M}\) and \(\mathcal{N}\) be two Banach manifolds. Let \(f \colon \mathcal{M} \to \mathcal{N}\) be smooth.
\begin{enumerate}
    \item We say that \(f\) is a \textit{(Banach)-immersion}\index{Banach!immersion} if \(\diff_mf\) is injective and its image is a closed subspace of \(T_{f(m)} \mathcal{N}\) that splits for each \(m\in \mathcal{M}\).
    \item We say that \(f\) is a \textit{(Banach)-submersion}\index{Banach!submersion} if \(\diff_m f\) is surjective and \(\ker \diff_mf\) splits for every \(m\in \mathcal{M}\). We say that \(f(m)\) is a \textit{regular value}\index{Regular value} of \(f\).
    \item We say that \(f\) is a \textit{(Banach)-embedding}\index{Banach!embedding} if it is an immersion and a homeomorphism into its image.
\end{enumerate}
}
If \(f\) is an embedding, then \(f(\mathcal{M}) \subset \mathcal{N}\) is a submanifold (in the sense defined above). It is sometimes called an \textit{embedded submanifold}. As we see, the main difference with respect to finite-dimensional manifolds is the split condition (which is always satisfied in finite-dimensional spaces). The split condition is needed to recover the local immersion and submersion theorems, that is, to see immersions as inclusions and submersions as projections in suitable coordinates. 

\subsection{Banach Lie groups}
Unsurprisingly, a \textit{Banach-Lie group}\index{Banach-Lie group} is defined as a group that has the structure of a Banach manifold, and such that the product and inversion maps are smooth, and a \textit{Hilbert-Lie group}\index{Hilbert-Lie group} is defined analogously. In this setting, right and left-translations \(L_g\) and \(R_g\) are diffeomorphisms, just as in the finite-dimensional case. The Lie algebra is also defined quite similarly to the finite-dimensional case. Firstly, two vector fields \(X \in \mathfrak{X}(\mathcal{M}), Y \in \mathfrak{X}(\mathcal{N})\) are \(f\)-\textit{related} (where \(f\colon \mathcal{M}\to \mathcal{N}\) is a smooth map) if \(\diff_mf(X_m)=Y_{f(m)}\) for all \(m \in \mathcal{M} \). We thus say that a vector field \(X \in \mathfrak{X}(G)\) is \textit{left-invariant} if it is \(L_g\)-related to itself. The \textit{Lie algebra} of \(G\), \(\mathrm{Lie}(G)\) is the set of all left invariant vector fields on \(G\), which can be identified with \(T_eG\) through the map \(X_p:=\diff_e L_g(X_e)\) (we will reserve the notation \(\mathfrak{g}\) for the Lie algebra of finite-dimensional Lie groups).

\Lemma{}{Let \(G\) be a Banach-Lie group. Then, \(\mathrm{Lie}(G)\) is a Lie subalgebra of \(\mathfrak{X}(G)\).}
 \begin{proof}
    It is obvious that \(\mathrm{Lie}(G)\) is a vector space, so we need only check that it is closed under the Lie bracket operation, but this follows from the fact that if \(X_i\) and \(Y_i\) are \(f\)-related, \(i=1,2\), then \([X_1, X_2]\) and \([Y_1, Y_2]\) are \(f\)-related as well (see \parencite[Page 117]{SergeLang}).
 \end{proof}

 Finally, Banach-Lie groups have a well-defined exponential map. This is not true for Lie groups modelled on arbitrary locally convex spaces. The definition is almost identical to the finite-dimensional one. If \(\alpha\colon I \to M\) is a smooth curve, we define 
 \begin{equation} \label{def_derivative_curve_Banach_manifold}
 \dot{\alpha}(t):=\diff_t\alpha\parent{\left. \frac{\diff}{\diff t} \right|_{t}} \in T_{\alpha(t)}\mathcal{M}.
 \end{equation}
 We say that a curve \(\alpha\) is an \textit{integral curve} of a vector field \(X\) if \(X_{\alpha(t)}=\dot{\alpha}(t)\) for every \(t \in I\). Given a Banach manifold, every vector field admits an integral curve defined on a maximal interval. The exponential map is just defined as
 \[
 \begin{array}{cccc}
    \exp \colon & \mathrm{Lie}(G) & \longrightarrow & G \\
    & X & \mapsto & \alpha(1)
 \end{array}
 \]
 where \(\alpha\) is the unique integral curve of \(X\) satisfying \(\alpha(0)=e\). See \parencite{Dodson_Galanis_Vassiliou_2015} for more details. Finally, the adjoint and coadjoint actions of \(G\) on \(\mathfrak{g}\) and \(\mathfrak{g}^*\) admit direct generalizations to this context. To define \(\Ad{}\), simply observe that the map \(C_g:=L_{g^{-1}} \circ R_g\) is a diffeomorphism, hence we may just define \(\Ad{g}(X):=\diff_e C_g(X)\) for \(X \in \mathfrak{g}\) and \(g \in G\). Similarly, we can just define \(\Ad{}^*\) by setting \((\Ad{g}^* \xi, X) :=(\xi, \Ad{g^{-1}}(X)).\)

\subsection{Differential forms in Banach manifolds.}
Banach manifolds are sufficiently well-behaved that we can still see differential forms as sections of adequate vector bundles (which is not the case for general infinite-dimensional manifolds, see for example \parencite[Appendix E]{Schmeding_2022} for details). The definition of vector bundle admits a straightforward generalization to Banach manifolds by simply asking them to be locally isomorphic to the product of the base manifold and the Banach space, and asking that, for every pair of trivializations \((U, \Phi), (V, \Psi)\), the map \(U \cap V\to \mathrm{Aut}(\mathbf{E}), m \mapsto \Phi_m \circ \Psi^{-1}_m \) is smooth.

%  Firstly, the definition of vector bundle can be adapted to a Banach-setting: A \textit{Banach vector bundle}\index{Vector bundle!Banach} over a Banach manifold is a Banach manifold \(E\), endowed with a smooth surjective map \(\pi\colon E \to \mathcal{M}\) satisfying the following properties:
% \begin{enumerate}
%     \item There is a Banach space \(\mathbf{E}\) such that each point \(x\in \mathcal{M}\) has an open neighborhood \(U\), and a diffeomorphism \(\Phi\colon \pi^{-1}(U) \to U_i \times \mathbf{E}\) that satisfies \(\pi=\mathrm{pr}_{U} \circ \Phi\).
%     \item Each fiber \(E_m:=\pi^{-1}(m)\) is a vector space, and \(\Phi_m:=\Phi|_{E_m} \colon E_m \to \corch{m} \times \mathbf{E}\) is a vector space isomorphism. The pairs \((U, \Phi)\) are called local trivializations.
%     \item For each pair of local trivializations \((U, \Phi)\), \((V, \Psi)\), \(\Phi_m \circ \Psi_m^{-1}\colon \mathbf{E} \to \mathbf{E}\) is a continuous vector space isomorphism whose inverse is continuous.
%     \item For each pair of local trivializations, (\(U, \Phi)\), \((V, \Psi)\), the map \(U \cap V \to \mathrm{Aut}(\mathbf{E}),m \mapsto (\Phi \circ \Psi^{-1})_m\) is smooth.
% \end{enumerate}

Let \(\mathrm{L}_{\mathrm{Alt}}^r(T\mathcal{M})\) be the vector bundle over \(\mathcal{M}\) whose fibers are the Banach spaces of multilinear alternating maps from \(T_m\mathcal{M}\) to \(\R\). A differential form on \(\mathcal{M}\) is defined as a section of this bundle. We keep the usual notation \(\Omega^p_\R(\mathcal{M})\) for \(p\)-forms. As in the finite-dimensional case, \(\Omega_\R^p(\mathcal{M})\) is a real vector space, and a \(\mathscr{C}^\infty(\mathcal{M};\R)\)-module. 

To define the exterior differential, one must observe that the usual definition in terms of local coordinates no longer makes sense in Banach manifolds. Instead, the exterior differential is defined using a coordinate-free relation that holds in the finite-dimensional setting as well (See \parencite[Prop. E.3]{Schmeding_2022} or \parencite[Prop. 3.2]{SergeLang}).

\Proposition{\textbf{Exterior differential}}{Let \(\mathcal{M}\) be a Banach manifold, and \(\omega \in \Omega^p_\R(\mathcal{M})\). Then, there is a unique \((p+1)\)-form on \(\mathcal{M}\) such that for all smooth vector fields \(X_1, \ldots X_{p+1}\in \mathfrak{X}(\mathcal{M})\), 
\begin{equation} \label{def_ext_differential}
    \begin{split}
        \diff\omega(X_1, \ldots ,X_{p+1})& = \sum_{i=1}^{p+1} (-1)^{i-1} X_i\big(\omega(X_1, \ldots, \hat{X}_i, \ldots, X_{p+1})\big)\\
        &+\sum_{1 \leq i<j \leq p+1}(-1)^{i+j} \omega([X_i, X_j], X_1, \ldots, \hat{X}_i, \ldots, \hat{X}_j, \ldots, X_{p+1}).
    \end{split}
\end{equation}
}
{\setlength{\parskip}{0pt}
% \begin{proof}
%     See \parencite[Prop. E.3]{Schmeding_2022} or \parencite[Prop. 3.2]{SergeLang}
% \end{proof}
\(\hspace{1 pt}\)\newline
Again, Banach manifolds are well-behaved enough so that most properties of the usual exterior differential are preserved. The exterior or wedge product is again defined analogously to the finite-dimensional case (see \parencite[Chapter V, section 3]{SergeLang} for details). }
\section{Basics of locally convex spaces} {\parskip=0pt
First, recall that a \textit{topological vector space}\index{Topological vector space}, TVS for short, is a vector space endowed with a topology such that addition and scalar products are continuous. Morphisms in the category \textbf{TVS} are continuous linear maps.

\Definition{\textbf{Seminorm}}{Let \(V\) be a TVS, a \textit{seminorm}\index{Seminorm} is a map \(p \colon V \to [0, \infty)\) satisfying
\begin{enumerate}
    \item \(p(\lambda v)=|\lambda| \cdot p(v)\) for every \(\lambda\in \R, v \in V\).
    \item \(p(v+w) \leq p(v)+p(w)\) for all \(v,w \in V\).
\end{enumerate}
}
Observe that a seminorm is just a norm for which the condition \(p(v)=0\) if and only if \(v=0\) is omitted. If we consider \(p_i^{-1}(0)\), which is a vector subspace of \(V\), then the quotient \(V/p_i^{-1}(0)\) has a natural norm \(||x+p_i^{-1}(0)||:=p(x)\) for any representative \(x\) of the equivalence class. Hence, the quotient has a natural topology as a normed vector space. We can now define locally convex spaces. 
\Definition{\textbf{Locally convex TVS}}{Let \(V\) be a TVS, we say that it is \textit{locally convex}\index{Locally convex space} if there is a family of continuous seminorms \(\corch{p_i}_{i \in I}\) such that:
\begin{enumerate}
    \item The topology of \(V\) is the initial topology with respect to the canonical projections \(q_i\colon V \to V/ p_i^{-1}(0)\).
    \item If \(v \in V\) and \(p_i(v)=0\) for all \(i \in I\), then \(v=0\).
\end{enumerate}
}
The second property of the definition implies that \(V\) is Hausdorff. \parencite[Prop. A.10]{Schmeding_2022} gives explicitly a subbase of the TVS topology on \(V\) and gives a proof of the Hausdorff condition. Observe that asking for \(V\) to be a TVS is superfluous. Indeed, if \(V\) is a vector space endowed with the initial topology with respect to the natural quotient maps then sum and product by scalars are continuous. Let us see it for sums. Let \(\tau\) be the initial topology on \(V\), and \(\tau_i\) be the metric topology on \(V_i:=V/p_i^{-1}(0)\) (which turns each \(V_i\) into a TVS). Consider the map \(s\colon V \times V \to V, (x,y) \mapsto x+y\). }

\noindent
\begin{minipage}{0.58\textwidth}
      By the universal property of the initial topology, \(s\) is continuous if and only if each \(q_i \circ s\) is continuous. However, observe that the diagram commutes. \(s_i\) is the sum in \(V_i\) (which is continuous). Hence, \(q_i \circ s=s_i \circ (q_i \times q_i)\), the right-hand side being continuous. This means that \(s\) is continuous. The same reasoning works for scalar multiplication.
\end{minipage} \hfill
\begin{minipage}{0.39\textwidth}
    \begin{center}
    \begin{tikzcd}[row sep=normal, column sep=normal, text height=1ex, text depth=1ex, ampersand replacement=\&]
    (V, \tau) \arrow[r, "{\displaystyle q_i}"] \& (V_i, \tau_i) \\
    (V \times V, \tau \times \tau) \arrow[r, " {\displaystyle q_i \times q_i}"] \arrow[u, "{\displaystyle s}", '] \& (V_i \times V_i, \tau_i \times \tau_i) \arrow[u, "{\displaystyle s_i}", ']
    \end{tikzcd}
\end{center}
\end{minipage}
% The second property of the definition implies that \(V\) is Hausdorff. \parencite[Prop. A.10]{Schmeding_2022} gives explicitly a subbase of the TVS topology on \(V\) and gives a proof of the Hausdorff condition. Observe that asking for \(V\) to be a TVS is superfluous. Indeed, if \(V\) is a vector space endowed with the initial topology with respect to the natural quotient maps then sum and product by scalars are continuous. Let us see it for sums. Let \(\tau\) be initial topology on \(V\), and \(\tau_i\) be the metric topology on \(V_i:=V/p_i^{-1}(0)\) (which turns each \(V_i\) into a TVS). Consider the map \(s\colon V \times V \to V, (x,y) \mapsto x+y\). By the universal property of the initial topology, \(s\) is continuous if and only if each \(q_i \circ s\) is continuous. However, observe that the diagram commutes. 
% \begin{center}
%     \begin{tikzcd}[row sep=normal, column sep=normal, text height=1ex, text depth=1ex, ampersand replacement=\&]
%     (V, \tau) \arrow[r, "{\displaystyle q_i}"] \& (V_i, \tau_i) \\
%     (V \times V, \tau \times \tau) \arrow[r, " {\displaystyle q_i \times q_i}"] \arrow[u, "{\displaystyle s}", '] \& (V_i \times V_i, \tau_i \times \tau_i) \arrow[u, "{\displaystyle s_i}", ']
%     \end{tikzcd}
% \end{center}
% where \(s_i\) is the sum in \(V_i\) (which is continuous). Hence, \(q_i \circ s=s_i \circ (q_i \times q_i)\), the right-hand side being continuous. This means that \(s\) is continuous. The same reasoning works for scalar multiplication. \\
\(\hspace{1 pt}\) \newline

Locally convex TVS are important because they retain one of the most celebrated and useful results of functional analysis: the Hahn-Banach theorem.

\Theorem{\textbf{Hahn-Banach}}{Let \(V\) be a locally convex TVS. Then, for every pair of distinct points \(x,y \in V\), there is a continuous linear map \(\alpha\colon V \to \R\) such that \(\alpha(x) \neq \alpha(y)\).}
% \begin{proof}
%     This is a standard result in functional analysis. See \parencite[Prop. 22.12]{meise1997introduction} for a proof.
% \end{proof}